% Borrador de sustitución del bloque FASE 4 de metodologia.tex.
% Se conservó la cronología descrita por el autor. Los detalles de software
% se contrastaron con pruebas de automatico v2 (consulta: 2026-09-07).
% No sustituir las Fases 1--3 ni las Fases 5--6.

En la cuarta fase se desarrollaron el accionamiento electrónico, el modo de operación manual y la referenciación del brazo. Después de comprobar el funcionamiento mecánico, se montó el circuito de prueba en placas de conexiones sin soldadura, o protoboards, y se utilizó cable calibre AWG~22 para las interconexiones entre módulos y microcontroladores.

La programación se inició con las funciones de movimiento de los motores mediante comandos. En la Portenta H7, mediante la interfaz Machine Control, se implementaron las señales de paso y dirección y el conteo de los pulsos generados. Posteriormente, se incorporó la lectura de los finales de carrera y el bloqueo del desplazamiento hacia un límite activo.

Se integró el mando inalámbrico mediante Bluepad32 en la ESP32. Sus entradas se transformaron en órdenes de dirección para los ejes X, Y y Z y se intercambiaron con la Portenta mediante I2C. La Portenta ejecutó las órdenes admitidas y devolvió información de movimiento, límites y estados. La ESP32 también gestionó los servomotores del efector final y una pantalla OLED como interfaz humano-máquina.

Para el intercambio entre placas, se configuró la Portenta como maestra I2C y la ESP32 como esclava. Se definieron mensajes binarios de longitud fija para transferir las órdenes y los estados, se incorporaron comprobaciones de integridad y vigencia y se programaron las respuestas ante fallos de comunicación. La OLED se conectó a un segundo bus I2C gestionado por la ESP32. El montaje se documentó mediante fotografías, en las cuales se identificó un controlador DM542T.

Después de comprobar esta interacción, se desarrolló la rutina de referenciación y determinación del recorrido en pasos. Para cada eje se programaron la aproximación a sus dos finales de carrera, la liberación de cada sensor y una segunda aproximación a menor frecuencia de pasos. A partir del conteo entre extremos se estableció una referencia central y se calcularon los factores de conversión de pasos por milímetro para X y Y.

Después del desarrollo del movimiento mediante joysticks, se incorporó el ingreso de coordenadas por terminal y una función de conversión de coordenadas cartesianas a destinos de pasos para el posicionamiento de X y Y. Esta función se distinguió del movimiento direccional mediante joystick, que utilizó órdenes de avance, retroceso y detención. Se emplearon recorridos de referencia de 446~mm para X y 336~mm para Y. Las comprobaciones funcionales descritas por el autor abarcaron el movimiento del brazo, la lectura de los finales y la interacción entre el mando, la ESP32 y la Portenta; no se les atribuyeron errores de posicionamiento ni tiempos de respuesta que no se hubieran documentado.

% PENDIENTES: alimentación, ajustes de corriente/micropasos y esquema físico.
% La foto identificó un DM542T, no todos los controladores ni su configuración.
% El autor confirmó Y = 336 mm como recorrido físico y confirmó la incorporación
% de coordenadas por terminal después del accionamiento mediante joysticks.
% El orden de desarrollo no equivalió al arranque actual: la versión V2
% también incluyó etapas de cámara y encoder, fuera del alcance de esta fase.
